% Options for packages loaded elsewhere
\PassOptionsToPackage{unicode}{hyperref}
\PassOptionsToPackage{hyphens}{url}
\documentclass[
]{article}
\usepackage{xcolor}
\usepackage[margin=1in]{geometry}
\usepackage{amsmath,amssymb}
\setcounter{secnumdepth}{5}
\usepackage{iftex}
\ifPDFTeX
  \usepackage[T1]{fontenc}
  \usepackage[utf8]{inputenc}
  \usepackage{textcomp} % provide euro and other symbols
\else % if luatex or xetex
  \usepackage{unicode-math} % this also loads fontspec
  \defaultfontfeatures{Scale=MatchLowercase}
  \defaultfontfeatures[\rmfamily]{Ligatures=TeX,Scale=1}
\fi
\usepackage{lmodern}
\ifPDFTeX\else
  % xetex/luatex font selection
\fi
% Use upquote if available, for straight quotes in verbatim environments
\IfFileExists{upquote.sty}{\usepackage{upquote}}{}
\IfFileExists{microtype.sty}{% use microtype if available
  \usepackage[]{microtype}
  \UseMicrotypeSet[protrusion]{basicmath} % disable protrusion for tt fonts
}{}
\makeatletter
\@ifundefined{KOMAClassName}{% if non-KOMA class
  \IfFileExists{parskip.sty}{%
    \usepackage{parskip}
  }{% else
    \setlength{\parindent}{0pt}
    \setlength{\parskip}{6pt plus 2pt minus 1pt}}
}{% if KOMA class
  \KOMAoptions{parskip=half}}
\makeatother
\usepackage{color}
\usepackage{fancyvrb}
\newcommand{\VerbBar}{|}
\newcommand{\VERB}{\Verb[commandchars=\\\{\}]}
\DefineVerbatimEnvironment{Highlighting}{Verbatim}{commandchars=\\\{\}}
% Add ',fontsize=\small' for more characters per line
\usepackage{framed}
\definecolor{shadecolor}{RGB}{248,248,248}
\newenvironment{Shaded}{\begin{snugshade}}{\end{snugshade}}
\newcommand{\AlertTok}[1]{\textcolor[rgb]{0.94,0.16,0.16}{#1}}
\newcommand{\AnnotationTok}[1]{\textcolor[rgb]{0.56,0.35,0.01}{\textbf{\textit{#1}}}}
\newcommand{\AttributeTok}[1]{\textcolor[rgb]{0.13,0.29,0.53}{#1}}
\newcommand{\BaseNTok}[1]{\textcolor[rgb]{0.00,0.00,0.81}{#1}}
\newcommand{\BuiltInTok}[1]{#1}
\newcommand{\CharTok}[1]{\textcolor[rgb]{0.31,0.60,0.02}{#1}}
\newcommand{\CommentTok}[1]{\textcolor[rgb]{0.56,0.35,0.01}{\textit{#1}}}
\newcommand{\CommentVarTok}[1]{\textcolor[rgb]{0.56,0.35,0.01}{\textbf{\textit{#1}}}}
\newcommand{\ConstantTok}[1]{\textcolor[rgb]{0.56,0.35,0.01}{#1}}
\newcommand{\ControlFlowTok}[1]{\textcolor[rgb]{0.13,0.29,0.53}{\textbf{#1}}}
\newcommand{\DataTypeTok}[1]{\textcolor[rgb]{0.13,0.29,0.53}{#1}}
\newcommand{\DecValTok}[1]{\textcolor[rgb]{0.00,0.00,0.81}{#1}}
\newcommand{\DocumentationTok}[1]{\textcolor[rgb]{0.56,0.35,0.01}{\textbf{\textit{#1}}}}
\newcommand{\ErrorTok}[1]{\textcolor[rgb]{0.64,0.00,0.00}{\textbf{#1}}}
\newcommand{\ExtensionTok}[1]{#1}
\newcommand{\FloatTok}[1]{\textcolor[rgb]{0.00,0.00,0.81}{#1}}
\newcommand{\FunctionTok}[1]{\textcolor[rgb]{0.13,0.29,0.53}{\textbf{#1}}}
\newcommand{\ImportTok}[1]{#1}
\newcommand{\InformationTok}[1]{\textcolor[rgb]{0.56,0.35,0.01}{\textbf{\textit{#1}}}}
\newcommand{\KeywordTok}[1]{\textcolor[rgb]{0.13,0.29,0.53}{\textbf{#1}}}
\newcommand{\NormalTok}[1]{#1}
\newcommand{\OperatorTok}[1]{\textcolor[rgb]{0.81,0.36,0.00}{\textbf{#1}}}
\newcommand{\OtherTok}[1]{\textcolor[rgb]{0.56,0.35,0.01}{#1}}
\newcommand{\PreprocessorTok}[1]{\textcolor[rgb]{0.56,0.35,0.01}{\textit{#1}}}
\newcommand{\RegionMarkerTok}[1]{#1}
\newcommand{\SpecialCharTok}[1]{\textcolor[rgb]{0.81,0.36,0.00}{\textbf{#1}}}
\newcommand{\SpecialStringTok}[1]{\textcolor[rgb]{0.31,0.60,0.02}{#1}}
\newcommand{\StringTok}[1]{\textcolor[rgb]{0.31,0.60,0.02}{#1}}
\newcommand{\VariableTok}[1]{\textcolor[rgb]{0.00,0.00,0.00}{#1}}
\newcommand{\VerbatimStringTok}[1]{\textcolor[rgb]{0.31,0.60,0.02}{#1}}
\newcommand{\WarningTok}[1]{\textcolor[rgb]{0.56,0.35,0.01}{\textbf{\textit{#1}}}}
\usepackage{longtable,booktabs,array}
\newcounter{none} % for unnumbered tables
\usepackage{calc} % for calculating minipage widths
% Correct order of tables after \paragraph or \subparagraph
\usepackage{etoolbox}
\makeatletter
\patchcmd\longtable{\par}{\if@noskipsec\mbox{}\fi\par}{}{}
\makeatother
% Allow footnotes in longtable head/foot
\IfFileExists{footnotehyper.sty}{\usepackage{footnotehyper}}{\usepackage{footnote}}
\makesavenoteenv{longtable}
\usepackage{graphicx}
\makeatletter
\newsavebox\pandoc@box
\newcommand*\pandocbounded[1]{% scales image to fit in text height/width
  \sbox\pandoc@box{#1}%
  \Gscale@div\@tempa{\textheight}{\dimexpr\ht\pandoc@box+\dp\pandoc@box\relax}%
  \Gscale@div\@tempb{\linewidth}{\wd\pandoc@box}%
  \ifdim\@tempb\p@<\@tempa\p@\let\@tempa\@tempb\fi% select the smaller of both
  \ifdim\@tempa\p@<\p@\scalebox{\@tempa}{\usebox\pandoc@box}%
  \else\usebox{\pandoc@box}%
  \fi%
}
% Set default figure placement to htbp
\def\fps@figure{htbp}
\makeatother
\setlength{\emergencystretch}{3em} % prevent overfull lines
\providecommand{\tightlist}{%
  \setlength{\itemsep}{0pt}\setlength{\parskip}{0pt}}
\usepackage{tcolorbox}
\usepackage{fancyvrb}
\usepackage{xcolor}
\usepackage{framed}
\definecolor{shadecolor}{HTML}{F7F7F7}
\usepackage{booktabs}
\usepackage{amsmath}
\usepackage{amssymb}
\usepackage{setspace}
\onehalfspacing
\usepackage{float}
\usepackage{bookmark}
\IfFileExists{xurl.sty}{\usepackage{xurl}}{} % add URL line breaks if available
\urlstyle{same}
% fallback for those not using the hyperref driver hyperxmp:
\makeatletter
\@ifundefined{xmpquote}{\newcommand{\xmpquote}[1]{#1}}{}
\makeatother
\hypersetup{
  pdftitle={someRpackages},
  pdfauthor={Simon Labarthe},
  hidelinks,
  pdfcreator={LaTeX via pandoc}}

\title{someRpackages}
\author{Simon Labarthe}
\date{August 2026}

\begin{document}
\maketitle

{
\setcounter{tocdepth}{2}
\tableofcontents
}
\section{Abstract}\label{abstract}

This article introduces \texttt{someRpackages}, a curated collection of
R functions designed to enhance efficiency in data analysis, automated
reporting, and AI-supported text classification. Researchers frequently
face repetitive tasks in data preparation and output generation. The
package addresses this by offering straightforward functions for
creating publication-ready LaTeX/PDF tables and integrating with AI APIs
for content coding. The core functions are presented with practical
examples obviously coming from \texttt{starwars}. By automating these
common analytical workflows, \texttt{someRpackages} allows researchers
to focus on substantive interpretation rather than technical
implementation. The package is open-source and available on GitHub.

\textbf{Keywords.} R package, Data Analysis, Automated Reporting, LaTeX,
Text Classification, API.

\newpage

\section{Introduction}\label{introduction}

Social scientists spend a substantial portion of their time not on
thinking, but on the mechanical aspects of data analysis: cleaning data,
writing repetitive code to generate tables, and manually classifying
text data. This is an inefficient use of intellectual resources (already
limited in my case). While R offers powerful analytical capabilities,
the path from a clean dataset to a polished, publication-ready table is
often fraught with cumbersome code and debugging.

\texttt{someRpackages} offers a pragmatic solution by wrapping complex R
operations into intuitive functions that produce LaTeX and PDF outputs.

\newpage

\section{Package Overview and Core
Functions}\label{package-overview-and-core-functions}

\texttt{someRpackages} is a collection of R utility functions designed
for data analysis and automated reporting. The package is designed to be
intuitive, requiring minimal R expertise to generate high-quality
outputs.

\subsection{Installation and Quick
Start}\label{installation-and-quick-start}

The package can be loaded directly from GitHub using the provided
script. To get started, run the following in your R console:

\begin{Shaded}
\begin{Highlighting}[]
\CommentTok{\# Load and install all functions with one line}
\FunctionTok{source}\NormalTok{(}\StringTok{"https://raw.githubusercontent.com/SimonLabarthe/someRpackages/main/load\_someRpackages.R"}\NormalTok{)}
\FunctionTok{load\_someRpackages}\NormalTok{()}
\end{Highlighting}
\end{Shaded}

This command sources and loads the core functions into the global
environment. For PDF and LaTeX output, a LaTeX distribution is required,
which can be conveniently installed via R:

\begin{Shaded}
\begin{Highlighting}[]
\FunctionTok{install.packages}\NormalTok{(}\StringTok{"tinytex"}\NormalTok{)}
\NormalTok{tinytex}\SpecialCharTok{::}\FunctionTok{install\_tinytex}\NormalTok{()}
\end{Highlighting}
\end{Shaded}

\subsection{Available Functions}\label{available-functions}

As of September 10, 2026, the package includes the following key
functions:

\begin{enumerate}
\def\labelenumi{\arabic{enumi}.}
\item
  \texttt{create\_descriptive\_table()}: Generates a comprehensive
  descriptive statistics table (including n, mean, median, mode, SD,
  min, max, and quartiles) for a given dataset. It creates a LaTeX/.tex
  file and compiles it into a standalone PDF, offering extensive
  customization options (e.g., selecting which statistics to display).
\item
  \texttt{create\_frequency\_table\_2vars()}: Creates a cross-tabulation
  frequency table for two variables, outputting results in a similar
  polished PDF format.
\item
  \texttt{create\_frequency\_table\_3vars()}: Extends the frequency
  table functionality to three variables, producing a nested frequency
  table that is particularly useful for subgroup analysis.
\item
  \texttt{categorise\_api()}: An advanced function that uses the
  DeepSeek API to classify a vector of texts. It allows researchers to
  efficiently code large text corpora based on a custom taxonomy, with
  built-in checkpointing and batch processing.
\end{enumerate}

\subsection{Design Philosophy}\label{design-philosophy}

The design philosophy behind \texttt{someRpackages} is guided by two key
principles. First, \textbf{laziness}: each function is designed to
perform a complete task with minimal user input. The goal is to reduce
complex data processing and output generation to a single readable line
of code. For example, generating a full descriptive statistics table in
PDF is as simple as:

\begin{Shaded}
\begin{Highlighting}[]
\FunctionTok{create\_descriptive\_table}\NormalTok{(data, }\AttributeTok{output\_dir =} \FunctionTok{getwd}\NormalTok{())}
\end{Highlighting}
\end{Shaded}

Second, \textbf{maximum coolness}: the emphasis on LaTeX and PDF
generation ensures that the output is ready for immediate integration
into manuscripts, reports, or presentations. This approach bridges the
gap between data analysis and scholarly publication.

\section{A Worked Example: Starwars}\label{a-worked-example-starwars}

To illustrate the package in action, we use the \texttt{starwars}
dataset bundled with \texttt{dplyr}.

\begin{Shaded}
\begin{Highlighting}[]
\FunctionTok{library}\NormalTok{(dplyr)}
\DocumentationTok{\#\# Warning: package \textquotesingle{}dplyr\textquotesingle{} was built under R version 4.5.3}
\DocumentationTok{\#\# }
\DocumentationTok{\#\# Attaching package: \textquotesingle{}dplyr\textquotesingle{}}
\DocumentationTok{\#\# The following objects are masked from \textquotesingle{}package:stats\textquotesingle{}:}
\DocumentationTok{\#\# }
\DocumentationTok{\#\#     filter, lag}
\DocumentationTok{\#\# The following objects are masked from \textquotesingle{}package:base\textquotesingle{}:}
\DocumentationTok{\#\# }
\DocumentationTok{\#\#     intersect, setdiff, setequal, union}
\FunctionTok{data}\NormalTok{(}\StringTok{"starwars"}\NormalTok{, }\AttributeTok{package =} \StringTok{"dplyr"}\NormalTok{)}
\end{Highlighting}
\end{Shaded}

\subsection{Descriptive Table}\label{descriptive-table}

\begin{Shaded}
\begin{Highlighting}[]
\FunctionTok{create\_descriptive\_table}\NormalTok{(}
\NormalTok{  starwars,}
  \AttributeTok{output\_name =} \StringTok{"starwars\_descriptives"}\NormalTok{,}
  \AttributeTok{output\_dir  =} \FunctionTok{getwd}\NormalTok{(),}
  \AttributeTok{mean        =} \ConstantTok{TRUE}\NormalTok{,}
  \AttributeTok{median      =} \ConstantTok{TRUE}\NormalTok{,}
  \AttributeTok{sd          =} \ConstantTok{TRUE}\NormalTok{,}
  \AttributeTok{min         =} \ConstantTok{TRUE}\NormalTok{,}
  \AttributeTok{max         =} \ConstantTok{TRUE}\NormalTok{,}
  \AttributeTok{mode        =} \ConstantTok{TRUE}
\NormalTok{)}
\end{Highlighting}
\end{Shaded}

\begin{Shaded}
\begin{Highlighting}[]
\NormalTok{starwars }\SpecialCharTok{\%\textgreater{}\%}
  \FunctionTok{select}\NormalTok{(height, mass, birth\_year) }\SpecialCharTok{\%\textgreater{}\%}
  \FunctionTok{summary}\NormalTok{()}
\DocumentationTok{\#\#      height           mass           birth\_year    }
\DocumentationTok{\#\#  Min.   : 66.0   Min.   :  15.00   Min.   :  8.00  }
\DocumentationTok{\#\#  1st Qu.:167.0   1st Qu.:  55.60   1st Qu.: 35.00  }
\DocumentationTok{\#\#  Median :180.0   Median :  79.00   Median : 52.00  }
\DocumentationTok{\#\#  Mean   :174.6   Mean   :  97.31   Mean   : 87.57  }
\DocumentationTok{\#\#  3rd Qu.:191.0   3rd Qu.:  84.50   3rd Qu.: 72.00  }
\DocumentationTok{\#\#  Max.   :264.0   Max.   :1358.00   Max.   :896.00  }
\DocumentationTok{\#\#  NA\textquotesingle{}s   :6       NA\textquotesingle{}s   :28        NA\textquotesingle{}s   :44}
\end{Highlighting}
\end{Shaded}

\subsection{Frequency Table}\label{frequency-table}

\begin{Shaded}
\begin{Highlighting}[]
\FunctionTok{create\_frequency\_table\_2vars}\NormalTok{(starwars, }\AttributeTok{var1 =} \StringTok{"gender"}\NormalTok{, }\AttributeTok{var2 =} \StringTok{"species"}\NormalTok{)}
\end{Highlighting}
\end{Shaded}

\begin{Shaded}
\begin{Highlighting}[]
\NormalTok{starwars }\SpecialCharTok{\%\textgreater{}\%}
  \FunctionTok{count}\NormalTok{(gender, species) }\SpecialCharTok{\%\textgreater{}\%}
  \FunctionTok{head}\NormalTok{(}\DecValTok{10}\NormalTok{)}
\end{Highlighting}
\end{Shaded}

{\def\LTcaptype{none} % do not increment counter
\begin{longtable}[]{@{}llr@{}}
\toprule\noalign{}
gender & species & n \\
\midrule\noalign{}
\endhead
\bottomrule\noalign{}
\endlastfoot
feminine & Clawdite & 1 \\
feminine & Droid & 1 \\
feminine & Human & 9 \\
feminine & Kaminoan & 1 \\
feminine & Mirialan & 2 \\
feminine & Tholothian & 1 \\
feminine & Togruta & 1 \\
feminine & Twi'lek & 1 \\
masculine & Aleena & 1 \\
masculine & Besalisk & 1 \\
\end{longtable}
}

\subsection{AI-Powered Text
Categorisation}\label{ai-powered-text-categorisation}

The \texttt{categorise\_api()} function classifies a vector of texts via
the DeepSeek API.

\begin{Shaded}
\begin{Highlighting}[]
\NormalTok{my\_texts }\OtherTok{\textless{}{-}} \FunctionTok{c}\NormalTok{(}
  \StringTok{"The new childcare policy will reduce costs for families."}\NormalTok{,}
  \StringTok{"We must fight for equal pay for equal work."}\NormalTok{,}
  \StringTok{"The traditional family unit is the foundation of our society."}
\NormalTok{)}

\NormalTok{api\_key }\OtherTok{\textless{}{-}} \FunctionTok{Sys.getenv}\NormalTok{(}\StringTok{"DEEPSEEK\_API\_KEY"}\NormalTok{)}

\NormalTok{results }\OtherTok{\textless{}{-}} \FunctionTok{categorise\_api}\NormalTok{(}
  \AttributeTok{texts           =}\NormalTok{ my\_texts,}
  \AttributeTok{api\_key         =}\NormalTok{ api\_key,}
  \AttributeTok{batch\_size      =} \DecValTok{2}\NormalTok{,}
  \AttributeTok{checkpoint\_file =} \StringTok{"my\_categorisation.rds"}
\NormalTok{)}

\FunctionTok{print}\NormalTok{(results)}
\end{Highlighting}
\end{Shaded}

The function returns a data frame with the same number of rows as the
input texts, containing the predicted \texttt{category} and
\texttt{theme} for each text.

\newpage

\section{Closing Statement}\label{closing-statement}

I believe that \texttt{someRpackages} can be a valuable asset for the
research community. I invite researchers to use it, contribute to it,
and adapt it to their needs. The source code is freely available on
GitHub, and I welcome any feedback or contributions :)

\subsubsection{Citation}\label{citation}

It took me a while to create all of this. I would really appreciate it
if you could make my work more visible.

You are pleased to cite as follow:

Labarthe, S. (2026). someRpackages: R Utility Functions.
\url{https://github.com/SimonLabarthe/someRpackages}

Visit My Website \url{https://github.com/SimonLabarthe}

\end{document}
